\documentclass[10pt,a4paper]{amsart}
\usepackage[margin=19mm]{geometry}
\usepackage{amsmath,amssymb,mathtools}
\usepackage[scr=rsfs,cal=euler]{mathalfa}
\usepackage{xfrac}
\usepackage[unicode,hidelinks,bookmarks=false]{hyperref}
\newcommand{\ie}{i.e.\ }
\newcommand{\cf}{cf.\ }
\newcommand{\resp}{respectively\ }
\newcommand{\Subsection}{\subsection}
\providecommand{\nobreakdash}{\nobreak\hbox{-}}
\setlength{\emergencystretch}{2em}
\multlinegap0pt
\usepackage{amssymb}
\usepackage{xcolor}
\usepackage{bm}
\newtheorem{lemma}{Lemma}
\newcommand{\sfx}{\mathsf{x}}
\title{\bf Phase transitions and approximations of mean squared error for state-space models with fractional differencing}
\author{Prabir Burman, Xiucai Ding,, Robert H. Shumway}
\date{Source: arXiv:2609.06727v1 (CC BY 4.0)}
\begin{document}
\maketitle

\noindent\emph{Source and licence.} \bf Phase transitions and approximations of mean squared error for state-space models with fractional differencing. Version arXiv:2609.06727v1 (CC BY 4.0), distributed by arXiv under the Creative Commons Attribution 4.0 International licence, http://creativecommons.org/licenses/by/4.0/, as recorded in the arXiv licence metadata for this exact version and shown on the abstract page https://arxiv.org/abs/2609.06727v1. Full licence notice: \texttt{licenses/arxiv-2609.06727-CC-BY-4.0.txt}.\par\smallskip

\noindent\textit{Editorial scope.} Counted unit: the article's lemma lem\_aux\_toeplitz (source0.tex lines 1811--1834). The article's complete original proof of it is delivered with it (source0.tex lines 1836--1931, one \texttt{proof} environment of 96 lines). No mathematics is written, repaired or re-derived: the statement and the proof are the article's own lines, in the article's own notation, and no retained line is substituted or re-flowed.  Retained uncounted as the source-local context the proof needs: lines 502--504.  Nothing was omitted from any retained span, so the package carries no omission record.  No retained line contains a citation, so the wrapper carries no bibliography and no bibliography entry.  No retained line contains a figure environment, an image-inclusion command or drawing code, so no image is delivered and the package declares no asset dependency.  Editorial formatting only, in addition: an AMS \texttt{amsart} preamble that restates the article's own declarations and macros used by the retained lines, namely the environments \texttt{lemma} on separate counters, together with the standard packages those lines need.  The retained environments print in this document as Lemma 1 in order of appearance, whereas the article prints the same items with the numbering of its own full document.  The complete native source of the article as distributed in the arXiv e-print for this exact version is bundled unchanged as \texttt{sources/209564/source0.tex}.
\begin{equation}\label{eq_points}
\sfx_j=\pi\frac{n+1-j}{n+1},
\end{equation}

\begin{lemma}\label{lem_aux_toeplitz}
Let \(T_n(g)\) be the Toeplitz matrix with symbol $g(u)=\sum_{j\in\mathbb{Z}}g_j\exp(-\mathrm{i}ju),$
where \(g_j=g_{-j}\) for every \(j\in\mathbb{Z}\) and $\sum_{j\in\mathbb{Z}}|jg_j|<\infty$. Define
\begin{equation*}
g_n(u)=\sum_{|j|\leq n}g_j\exp(-\mathrm{i}ju),
\end{equation*}
and let \(R_n(g)=(R_n(j,k;g))_{j,k=1}^n\), where we recall (\ref{eq_points}) and denote
\begin{align*}
R_n(j,k;g)
&=
\frac{2}{n+1}
\sum_{t=1}^{n}
g(\sfx_t)\sin(j\sfx_t)\sin(k\sfx_t), \ \,
1\leq j,k\leq n. 
\end{align*}
Then we have that 
\begin{equation*}
\left\|
T_n(g)-R_n(g)-H_n(g)-\mathcal{W}_nH_n(g)\mathcal{W}_n
\right\|
=
\mathrm{O}\!\left(\left\|g_n-g\right\|_\infty\right).
\end{equation*}
\end{lemma}

\begin{proof}
%Let \(U_n\) denote the discrete sine-transform matrix with \((j,t)\)th entry \(\{2/(n+1)\}^{1/2}\sin(j\sfx_t)\). The matrix \(U_n\) is orthogonal. For the trigonometric polynomial \(g_n\), the discrete sine-transform identity converts the sine-weighted integral below into the corresponding finite sum over \(\sfx_1,\ldots,\sfx_n\).
%
%
%Let \(U_n\) denote the discrete sine-transform matrix introduced above. Its \((j,t)\)th entry is \(\{2/(n+1)\}^{1/2}\sin(j\sfx_t)\). The discrete sine functions are orthogonal on the grid \(\{\sfx_t\}_{t=1}^n\), and hence \(U_n\) is orthogonal. Moreover, by the definition of \(R_n(g_n)\), its \((j,k)\)th entry is the \((j,k)\)th entry of
%\(U_n\operatorname{diag}\{g_n(\sfx_1),\ldots,g_n(\sfx_n)\}U_n^\top\). Thus, the finite sum below is precisely \(R_n(j,k;g_n)\).
First, suppose that \(j+k\leq n\). By the definitions of the Toeplitz and Hankel entries and the identity \(\cos((j-k)u)-\cos((j+k)u)=2\sin(ju)\sin(ku)\), we have
\begin{align*}
T_n(j-k;g_n)-H_n(j+k;g_n)
&=
\frac{2}{\pi}
\int_0^\pi
\sin(ju)\sin(ku)g_n(u)\,\mathrm{d}u \\
&=
\frac{2}{n+1}
\sum_{t=1}^{n}
\sin(j\sfx_t)\sin(k\sfx_t)g_n(\sfx_t) \\
&=
R_n(j,k;g_n).
\end{align*}
Indeed, the second equality is the discrete sine-transform quadrature formula. Since \(g_n\) is an even trigonometric polynomial of degree at most \(n\), this formula follows by expanding \(g_n\) in its Fourier series and applying discrete cosine orthogonality on the grid \(\{\sfx_t\}_{t=1}^n\). The condition \(j+k\leq n\) ensures that no aliasing terms arise in this calculation.




Next, suppose that \(n+1\leq j+k\leq 2n\). Applying the same discrete sine-transform identity to the indices \(n+1-j\) and \(n+1-k\) gives
\begin{align*}
&T_n(j-k;g_n)-H_n(2n+2-j-k;g_n) \\
&\qquad =
\frac{2}{\pi}
\int_0^\pi
\sin ((n+1-j)u)
\sin ((n+1-k)u)
g_n(u)\,\mathrm{d}u \\
&\qquad =
\frac{2}{n+1}
\sum_{t=1}^{n}
\sin ((n+1-j)\sfx_t)
\sin ((n+1-k)\sfx_t)
g_n(\sfx_t) \\
&\qquad =
\frac{2}{n+1}
\sum_{t=1}^{n}
\sin(j\sfx_t)\sin(k\sfx_t)g_n(\sfx_t) \\
&\qquad =
R_n(j,k;g_n).
\end{align*}
For the penultimate equality, note that the two sine factors acquire the same sign under the replacement \(j\mapsto n+1-j\) and \(k\mapsto n+1-k\), so their product is unchanged.

Because \(g_n\) has Fourier coefficients equal to zero outside \(\{-n,\ldots,n\}\), \(H_n(j+k;g_n)=0\) whenever \(j+k>n\). Moreover, the \((j,k)\)th entry of \(\mathcal{W}_nH_n(g_n)\mathcal{W}_n\) is \(H_n(2n+2-j-k;g_n)\), which is zero whenever \(j+k\leq n+1\). Thus, the first identity applies when \(j+k\leq n\), both Hankel terms vanish when \(j+k=n+1\), and the second identity applies when \(j+k\geq n+2\). Combining these cases entrywise yields
\begin{equation*}
T_n(g_n)-R_n(g_n)
=
H_n(g_n)+\mathcal{W}_nH_n(g_n)\mathcal{W}_n.
\end{equation*}

It remains to replace \(g_n\) by \(g\). Put \(h=g-g_n\). For the Hankel term, \(H_n(h)\) is a finite compression of the Hankel operator associated with the symbol \(h\). This operator is obtained by composing multiplication by \(h\), which has operator norm \(\|h\|_\infty\), with an orthogonal projection and a reflection, both of which have operator norm one. Therefore,
\begin{equation*}
\left\|H_n(g)-H_n(g_n)\right\|
=
\left\|H_n(h)\right\|
\leq
\left\|h\right\|_\infty
=
\left\|g-g_n\right\|_\infty.
\end{equation*}

For the sine-basis term, let \(U_n\) be the orthogonal discrete sine-transform matrix with \((j,t)\)th entry \((2/(n+1))^{1/2}\sin(j\sfx_t)\). By the definition of \(R_n(\cdot)\), \(R_n(h)\) is obtained by conjugating the diagonal matrix with diagonal entries \(h(\sfx_1),\ldots,h(\sfx_n)\) by \(U_n\). Orthogonal conjugation preserves the operator norm. Hence,
\begin{equation*}
\left\|R_n(g)-R_n(g_n)\right\|
=
\left\|R_n(h)\right\|
\leq
\max_{1\leq t\leq n}|h(\sfx_t)|
\leq
\left\|g-g_n\right\|_\infty.
\end{equation*}

Moreover, \(\mathcal{W}_n\) is an orthogonal and symmetric matrix; that is, \(\mathcal{W}_n^\top=\mathcal{W}_n=\mathcal{W}_n^{-1}\). Therefore, conjugation by \(\mathcal{W}_n\) preserves the operator norm, and hence
\begin{equation*}
\left\|
\mathcal{W}_n
\bigl (H_n(g)-H_n(g_n)\bigr)
\mathcal{W}_n
\right\|
=
\left\|
H_n(g)-H_n(g_n)
\right\|.
\end{equation*}
Furthermore, the \((j,k)\)th entry of \(T_n(g)\) is \(g_{j-k}\), where \(1\leq j,k\leq n\). Thus, only Fourier coefficients \(g_\ell\) with \(|\ell|\leq n-1\) enter \(T_n(g)\). Since \(g_n\) retains all Fourier coefficients of \(g\) with \(|\ell|\leq n\), the matrices \(T_n(g)\) and \(T_n(g_n)\) have identical entries. Therefore,
\begin{equation*}
T_n(g)=T_n(g_n).
\end{equation*}
Combining this equality with the identity established above for \(g_n\), the difference in the statement of the lemma is the sum of the differences between the corresponding \(R_n\) and \(H_n\) terms. The triangle inequality, together with the preceding bounds, shows that its operator norm is bounded by \(3\|g-g_n\|_\infty\). This proves the result.
\end{proof}

\end{document}
